\documentclass[11pt,a4paper]{article}

% ---- Encodage & langue ----
\usepackage[utf8]{inputenc}
\usepackage[T1]{fontenc}
\usepackage[french]{babel}

% ---- Mise en page ----
\usepackage[margin=2.2cm]{geometry}
\usepackage{titlesec}
\usepackage{enumitem}
\usepackage{xcolor}
\usepackage{array}
\usepackage{booktabs}
\usepackage{fancyhdr}

% ---- Couleurs sobres ----
\definecolor{primary}{HTML}{2C5F8A}
\definecolor{lightgray}{HTML}{F2F4F7}

% ---- En-tête / pied de page ----
\pagestyle{fancy}
\fancyhf{}
\lhead{\small\textcolor{primary}{Ausculta}}
\rhead{\small Fiche technique}
\cfoot{\small\thepage}
\renewcommand{\headrulewidth}{0.4pt}

% ---- Style des titres ----
\titleformat{\section}
  {\normalfont\large\bfseries\color{primary}}{\thesection.}{0.6em}{}
\titleformat{\subsection}
  {\normalfont\normalsize\bfseries}{\thesubsection}{0.6em}{}

\setlist[itemize]{leftmargin=1.4em,itemsep=2pt,topsep=3pt}

\begin{document}

% ================= EN-TÊTE =================
\begin{center}
  {\Huge\bfseries\color{primary} Ausculta}\\[4pt]
  {\large Logiciel de gestion de cabinet médical}\\[2pt]
  \textit{Fiche technique}\\[6pt]
  \small Version 0.0.0 \quad\textbullet\quad \today
\end{center}

\vspace{4pt}
\hrule
\vspace{10pt}

% ================= PRÉSENTATION =================
\section{Présentation}

\textbf{Ausculta} est une application de bureau destinée aux médecins
exerçant en cabinet privé. Elle centralise la gestion des patients, des
rendez-vous, des ordonnances, des documents médicaux et des statistiques
d'activité au sein d'une interface unique.

\medskip
L'application est \textbf{100\% hors-ligne} : toutes les données sont
stockées localement sur le poste du praticien, sans serveur distant ni
connexion Internet requise. La confidentialité des données médicales est
ainsi garantie par conception.

% ================= FONCTIONNALITÉS =================
\section{Fonctionnalités principales}

\subsection*{Gestion des patients}
\begin{itemize}
  \item Création, consultation, modification et suppression des fiches patients.
  \item Recherche rapide par nom ou par critère.
  \item Fiche détaillée regroupant informations personnelles, historique
        des rendez-vous, ordonnances et documents.
\end{itemize}

\subsection*{Rendez-vous}
\begin{itemize}
  \item Prise de rendez-vous avec date, heure, durée et motif de consultation.
  \item Agenda journalier et vue par plage de dates.
  \item Gestion des statuts (planifié, annulé, honoré, absence).
  \item Détection automatique des rendez-vous manqués au démarrage.
\end{itemize}

\subsection*{Ordonnances}
\begin{itemize}
  \item Rédaction d'ordonnances multi-médicaments (dosage, fréquence,
        quantité, durée, notes).
  \item Génération automatique d'ordonnances au format PDF à partir d'un
        modèle personnalisé.
  \item Historique complet des prescriptions par patient.
\end{itemize}

\subsection*{Documents médicaux}
\begin{itemize}
  \item Importation et archivage des documents liés à chaque patient.
  \item Ouverture directe des fichiers depuis l'application.
  \item Rattachement automatique des ordonnances générées à la fiche patient.
\end{itemize}

\subsection*{Statistiques \& tableau de bord}
\begin{itemize}
  \item Statistiques financières et volume de consultations sur une période.
  \item Taux d'absentéisme (\textit{no-show}) et suivi de l'activité.
  \item Tableau de bord synthétique et graphiques interactifs.
\end{itemize}

\subsection*{Profil \& paramètres}
\begin{itemize}
  \item Profil du médecin (nom, spécialité, coordonnées) utilisé sur les
        ordonnances.
  \item Interface multilingue : français, anglais et arabe (support
        d'affichage de droite à gauche).
  \item Authentification sécurisée et session persistante optionnelle.
\end{itemize}

% ================= CE QUE L'APP APPORTE =================
\section{Ce que l'application apporte}
\begin{itemize}
  \item \textbf{Gain de temps} : centralisation de toutes les tâches
        administratives du cabinet.
  \item \textbf{Sécurité} : données médicales stockées et chiffrées
        localement, sans exposition réseau.
  \item \textbf{Simplicité} : interface claire, sans configuration
        technique ni maintenance de serveur.
  \item \textbf{Accessibilité} : disponible en trois langues et adaptée aux
        besoins des cabinets francophones et arabophones.
\end{itemize}

% ================= FICHE TECHNIQUE =================
\section{Caractéristiques techniques}

\renewcommand{\arraystretch}{1.4}
\begin{center}
\begin{tabular}{>{\raggedright}p{4.5cm} p{9.5cm}}
\toprule
\textbf{Type d'application} & Logiciel de bureau (Windows, macOS, Linux) \\
\textbf{Technologies} & Electron, React, TypeScript \\
\textbf{Base de données} & SQLite locale (hors-ligne) \\
\textbf{Sécurité} & Mots de passe chiffrés (bcrypt), jetons JWT, stockage sécurisé \\
\textbf{Génération PDF} & Ordonnances via modèle personnalisable \\
\textbf{Langues} & Français, Anglais, Arabe (RTL) \\
\textbf{Installation} & Installateur Windows (NSIS), image macOS (dmg), AppImage (Linux) \\
\bottomrule
\end{tabular}
\end{center}

\vspace{12pt}
\begin{center}
\small\textcolor{primary}{Ausculta — Simplifiez la gestion de votre cabinet médical.}
\end{center}

\end{document}
